\chapter{Shipping Indradhanu V2}\label{ch:shipping}
\begingroup\sloppy % long provenance tags need loose lines rather than overfull ones

\begin{onething}
Putting FT04 into the product took six small decisions, each forced by something the product already did: a new
codename because codename routing is last-wins; a manifest that carries the whole model configuration because the api
cannot load a checkpoint; a band-grid check because one band filter feeds every model in an action;
\texttt{read\_bands} because the model must be told which wavelengths it reads; cloud zeroed \emph{after} scoring
because that is the order the benchmark measured; and a parity run that proves the product path gives the evaluator's
numbers.
\end{onething}

\section{From a checkpoint to a product}\label{sec:shipping:why}

\Cref{ch:ftruns} ended with a release candidate: FT04, scored with residual Mahalanobis, with cloud zeroed on its own
masks. On the real PRISMA bench that reads pAUC 0.170. But that number came from the \emph{evaluator}
(\code{app/evaluation/}, \cref{ch:runeval}), a research harness that loads checkpoints directly, on a laptop. The
\emph{product} (\cref{ch:product}) is something else: an api that must stay light and never imports torch, a worker
that runs actions from a Postgres queue, a chain of actions in which one \code{band_filter_apply} output feeds an
\code{anomaly_scoring} action, and defaults that every analyst inherits without choosing them.

The obvious way to ship a model - copy the checkpoint into the models folder and point the old code at it - fails in
six different places, and the six sections below are those six places. On 2026-10-19 FT04 was served in
\code{anomaly_scoring} under the codename \textbf{Indradhanu V2}, opt-in next to the shipped v0.2.0 ``Indradhanu'',
which stayed the default. The design was agreed first (\file{docs/proposals/ship-ft04.md}) in two phases: Phase 1 makes
V2 selectable without changing any default; Phase 2, later, would make it the hyperspectral default.

\begin{analogy}
A hospital takes on a new specialist who trained abroad. Before her first patient, the hospital gives her a name badge
that no other doctor on the ward shares (or the switchboard puts calls through to whoever was added last), files her
credentials in a form the front desk can read without phoning the overseas university, checks every referral is in the
units she works in, and has her see a few patients the old consultant already examined - to confirm she reaches the
same findings by the same route before anyone relies on her. Each of those four steps is a section of this chapter.
\end{analogy}

\begin{table}[htbp]
\centering
\caption{What was shipped on 2026-10-19, from the manifest \texttt{current.json} on the dev laptop.}
\label{tab:shipping:manifest}
\small
\begin{tabular}{ll}
\toprule
field & value \\
\midrule
architecture & \texttt{spectral\_stem\_finetune} \\
codename & Indradhanu V2 \\
version, epoch & \texttt{v2.0.0}, epoch 10 (of fine-tuning) \\
params & 5,886,451 \\
val\_loss & 0.05578 \\
checkpoint file & \texttt{spectral\_stem\_finetune\_v2.0.0\_epoch10.pt} \\
\bottomrule
\end{tabular}

\smallskip
{\footnotesize Source: \tagM{C:/es/ship/spectral_stem_finetune/current.json}.}
\end{table}

\begin{drift}
Since 2026-10-29 the codename Indradhanu V2 points to Indradhanu Heavy (FTH3 epoch 5, version \texttt{v2.1.0}, 35.3\,M
parameters; \cref{ch:heavy}), behind the same architecture key, with FT04 (\texttt{v2.0.0}) kept as the manifest's
first alternative (\file{docs/lld/spectral-stem.md}, ``Serving''). Everything in this chapter describes the FT04 ship
and still holds for the serving path, except where a box says otherwise.
\end{drift}

\section{Why the codename had to change}\label{sec:shipping:codename}

\index{codename routing}
A \emph{codename} is the name an analyst picks in \code{anomaly_scoring}: ``Indradhanu'', ``MNF-RX''. The worker turns
it into a model with one dictionary:

\begin{lstlisting}
# _anomaly_scoring_run.py (and anomaly_scoring.py's validate_config)
by_codename = {m.codename.lower(): m for m in catalog}
\end{lstlisting}

\code{catalog} comes from \code{list_catalog} in \code{backend/allotrope/foundation_models/resolver.py}, which walks
the model directories in sorted order (\code{for p in sorted(models_root.iterdir())}). Two facts decide which model a
repeated codename reaches.

\paragraph{Fact 1: a dictionary built this way keeps the last value.} The comprehension assigns the entries one at a
time, in catalog order. Assigning to a key that already exists overwrites its value. So if two manifests both say
``Indradhanu'', the dictionary ends up holding whichever came \emph{last}.

\paragraph{Fact 2: strings sort by character codes, left to right.} \index{lexicographic order}
Python compares two strings \emph{lexicographically}: compare the first characters; if they differ, the string with
the smaller character code comes first; if they are equal, move to the second characters, and so on. Character codes
are the Unicode code points: \texttt{a} is 97, \texttt{h} is 104, \texttt{s} is 115.

\begin{worked}[label={we:shipping:sort}]{Which manifest wins ``Indradhanu''?}
The two directories are \texttt{hyperspectral\_segformer\_mae} (v0.2.0) and \texttt{spectral\_stem\_finetune} (FT04).
Compare first characters: \texttt{h} is 104, \texttt{s} is 115.
\begin{align*}
104 &< 115\\
\Rightarrow\ &\texttt{hyperspectral\_segformer\_mae} \text{ sorts first, } \texttt{spectral\_stem\_finetune} \text{ last}\\
\Rightarrow\ &\text{``indradhanu''} \to \text{v0.2.0, then overwritten by FT04: FT04 wins.}
\end{align*}
The script confirms both the sort and the surviving value\tagC.
\end{worked}

Every default hyperspectral action (\code{_DEFAULT_PRISMA}, \code{_DEFAULT_ENMAP} and \code{_DEFAULT_AVIRIS_NG} in
\code{anomaly_scoring.py} all say \code{"model_codenames": ["Indradhanu"]}) would have been re-routed to FT04 without
anyone choosing it - and, with the grid check of \cref{sec:shipping:grid}, every one of them would then fail, because
the default band filter resamples to the 165-band grid. \code{list_catalog} does log a warning on a collision
(``codename collision \ldots\ Rename one of them''), but a log line is not a guard. Hence a new name.

\begin{keypoint}
When routing depends on the order of a directory listing, the only safe collision is no collision: give every model a
name no other model shares.
\end{keypoint}

\section{A manifest that carries the whole configuration}\label{sec:shipping:manifest}

\index{manifest}
The api lists models and validates actions; it never loads torch (the repo's lazy-import rule: action-type modules import
torch only inside the worker's \code{run()}). So everything the resolver needs must be readable from \file{current.json},
a JSON file next to each checkpoint. For v0.2.0, \code{_model_config_from_manifest} rebuilds the model configuration from
a few summary fields. For the spectral-stem models that is not possible:

\begin{lstlisting}
# backend/allotrope/foundation_models/resolver.py, _model_config_from_manifest
if arch in ("spectral_stem_finetune", "spectral_stem_mae"):
    # The stem configs carry far more than a manifest summary can rebuild (sources,
    # band weighting, the stem variant), and the api cannot torch.load; the manifest
    # therefore embeds the checkpoint's full config.model_config.
    cfg = current.get("model_config")
    if not isinstance(cfg, dict) or cfg.get("model_type") != arch:
        raise KeyError(f"{arch}: current.json must embed the checkpoint's model_config")
    return cfg
\end{lstlisting}

\code{scripts/make_model_manifest.py} is the one place that does load the checkpoint: it copies
\code{ck["config"]["model_config"]} into \code{current.model_config}, counts the parameters and copies the checkpoint
next to it as \texttt{<architecture>\_<version>\_epoch<N>.pt}. If the embedded config is missing or names another
architecture, the resolver raises \code{KeyError} - and \code{list_catalog} skips manifests that raise \code{KeyError},
so a bad manifest makes the model \emph{disappear} from the picker rather than load wrongly. Failing closed is the right
direction: a missing model is noticed; a wrongly configured one quietly scores.

The capabilities entry sets \code{pixel_stats_relpath=None}: there is no per-band statistics file, because the
normalisation curve travels inside the checkpoint as buffers (\texttt{stats.wl}, \texttt{stats.mean},
\texttt{stats.std}; \cref{ch:finetune}).

\subsection{What ``params'' counts}

\index{parameter!versus buffer}
A model's \code{state_dict} holds two kinds of tensors. \emph{Parameters} are the numbers training changes by gradient
descent. \emph{Buffers} are stored numbers that are not trained by gradients - BatchNorm's running statistics, or here
the normalisation curve and the trained-wavelength record. \code{make_model_manifest.py} counts
\begin{lstlisting}
"params": int(sum(v.numel() for k, v in sd.items()
                  if not k.startswith(("stats.", "trained_wl"))))
\end{lstlisting}
that is, every tensor in the state dict except the normalisation and wavelength buffers. Predict the figure from the
model's parts.

\begin{worked}[label={we:shipping:params}]{Where 5,886,451 comes from}
Trainable parameters of the shipped FT04, by top-level module (from the checkpoint\tagC):
\begin{align*}
\text{spectral stem in (\texttt{spectral\_in})} &\quad 14{,}721\\
\text{spectral decoder out (\texttt{spectral\_out})} &\quad 4{,}257\\
\text{BatchNorm scale and shift (\texttt{norm}, } 2 \times 32) &\quad 64\\
\text{encoder} &\quad 3{,}331{,}072\\
\text{decoder} &\quad 2{,}164{,}736\\
\text{segmentation head (\cref{we:ftruns:seghead})} &\quad 371{,}536
\end{align*}
Add them: $14{,}721 + 4{,}257 = 18{,}978$; $18{,}978 + 64 = 19{,}042$; $19{,}042 + 3{,}331{,}072 = 3{,}350{,}114$;
$3{,}350{,}114 + 2{,}164{,}736 = 5{,}514{,}850$; $5{,}514{,}850 + 371{,}536 = 5{,}886{,}386$ parameters.

The manifest says 5,886,451. The difference, $5{,}886{,}451 - 5{,}886{,}386 = 65$, is the BatchNorm's buffers, which the
manifest's filter does not exclude: running mean (32 numbers), running variance (32) and the step counter
\texttt{num\_batches\_tracked} (1): $32 + 32 + 1 = 65$.
\end{worked}

\begin{sourcenote}[: the manifest's ``params'' includes 65 buffers]
\code{scripts/make_model_manifest.py} (line 59) excludes only \texttt{stats.*} and \texttt{trained\_wl*} from its count,
so the reported \texttt{params} of 5,886,451 includes the stem BatchNorm's 65 buffer entries. Building the model and
summing \code{model.parameters()} gives 5,886,386 (\file{research/textbook/tools/compute_shipping.py}). The course
quotes the manifest's 5,886,451 (``about 5.9\,M''); both round to that. The excluded buffers are \texttt{trained\_wl}
(733 wavelengths in the checkpoint - the record of \cref{ch:ftruns}) and three 202-point normalisation curves.
\end{sourcenote}

\section{One band filter, two kinds of model: the grid check}\label{sec:shipping:grid}

\index{band grid check}
v0.2.0 reads the 165-band common grid (\cref{ch:sensors}). V2 reads each sensor's \emph{native} bands - the wavelengths
the instrument actually measured. \code{band_filter_apply} can produce either (\code{use_common_wavelength_grid},
default \texttt{true}), but an \code{anomaly_scoring} action reads \emph{one} upstream band-filter output for every
model it runs. So the resolver now records which grid each model needs:

\begin{lstlisting}
# resolver.py, ModelCapabilities
band_grid: str = "common165"          # v0.2.0 and every older model
# the spectral_stem_finetune entry:
band_grid="native", exclude_classes=("cloud",)
\end{lstlisting}

and the worker checks the vendable before running any foundation model on a hyperspectral scene:

\begin{lstlisting}
# _anomaly_scoring_run.py, _check_band_grid
grid = np.asarray(DEFAULT_COMMON_WAVELENGTH_GRID, dtype=np.float64)     # (165,)
on_grid = len(wavelengths) == len(grid) and np.allclose(wavelengths, grid)
if band_grid == "native" and on_grid:
    raise ValueError("... reads native bands: run band_filter_apply with "
                     "use_common_wavelength_grid=false ...")
if band_grid == "common165" and not on_grid:
    raise ValueError("... needs the 165-band grid: run band_filter_apply with "
                     "use_common_wavelength_grid=true ...")
\end{lstlisting}

\paragraph{What \texttt{np.allclose} tests.} \index{np.allclose}
\code{np.allclose(a, b)} is true when every pair of entries satisfies
\begin{equation}
\abs{a_i - b_i} \le \mathrm{atol} + \mathrm{rtol} \cdot \abs{b_i}, \qquad \mathrm{rtol} = 10^{-5},\ \mathrm{atol} =
10^{-8} \ \text{(the defaults)}. \label{eq:shipping:allclose}
\end{equation}
The tolerance grows with the size of the number: a relative tolerance.

\begin{worked}[label={we:shipping:allclose}]{How close is ``on the grid''?}
At the grid's first wavelength, 460\,nm, the allowed difference is
\[
10^{-8} + 10^{-5} \times 460 = 0.00000001 + 0.0046 = 0.0046\,\text{nm}.
\]
At its last, 2450\,nm: $10^{-8} + 10^{-5} \times 2450 = 0.0245$\,nm. So a list shifted by 0.001\,nm everywhere still
counts as the grid (0.001 is under 0.0046), while a shift of 0.01\,nm does not (0.01 is over 0.0046 at the short end).
The script agrees\tagC. No real sensor lists its native bands within a few thousandths of a nanometre of a 10\,nm grid,
so the test separates ``resampled'' from ``native'' with a wide margin.
\end{worked}

The test is ``do these wavelengths equal the grid'', not ``are there 165 of them''. A native list that happens to have
165 bands still fails \code{np.allclose} and counts as native. Classical detectors (MNF-RX and friends) skip the check:
they take either grid.

\begin{sourcenote}[: the check runs in the worker, not at submit time]
The design (\file{docs/proposals/ship-ft04.md}, Phase 1 step 4) placed the grid check in
\code{anomaly_scoring.validate_config}. The code runs it in the worker (\code{_check_band_grid},
\code{backend/allotrope/action_types/_anomaly_scoring_run.py}, line 93, called at line 297), because the api cannot see
the upstream band filter's settings when the action is submitted (\file{docs/09-known-issues.md}). The consequence:
\textbf{one action cannot run V2 and v0.2.0 together} - whichever does not match the band filter makes the action fail at
run time.
\end{sourcenote}

\section{\texttt{read\_bands}: tell the model what it is reading}\label{sec:shipping:readbands}

\index{trained band mask}
A spectral stem reads a band list of any length, but only if it is told the wavelengths (\cref{ch:crossstem}). The
product's older path never called \code{set_wavelengths}. The new helper does that and drops the bands the model has
never trained on:

\begin{lstlisting}
# app/foundation_models/inferencers/spectral_stem_inferencer.py
def read_bands(self, cube, wavelengths_nm):          # cube: (N, H, W)
    wl = np.asarray(wavelengths_nm, dtype=np.float64)
    keep = self.model.inner.trained_band_mask(wl).cpu().numpy()   # (N,) bool
    self.set_wavelengths(wl[keep])
    return cube[keep], wl[keep], np.flatnonzero(keep)  # (N_kept, H, W), (N_kept,), indices
\end{lstlisting}

\code{trained_band_mask} keeps a band when its distance to the nearest recorded wavelength is at most 15\,nm:
\begin{equation}
\text{keep}(\lambda) = \Big[\min_{r \in R} \abs{\lambda - r} \le 15\,\text{nm}\Big], \label{eq:shipping:keep}
\end{equation}
with $R$ the trained-wavelength record (\code{trained_wl}), which grew during fine-tuning as FT04 read EnMAP, PRISMA and
AVIRIS-NG. In code it is one line: \code{(wl[:, None] - self.trained_wl[None, :]).abs().min(dim=1).values <= tol_nm} -
a matrix of all band-to-record distances, shape \shape{(N, |R|)}, then the minimum along each row.

\begin{worked}[label={we:shipping:keep}]{Three bands against a record ending 2395, 2405, 2440}
Build the distance matrix row by row (do it on paper).
\begin{center}
\begin{tabular}{lccccc}
\toprule
band & to 2395 & to 2405 & to 2440 & minimum & keep? \\
\midrule
2400\,nm & 5 & 5 & 40 & 5 & $5 \le 15$: keep \\
2452\,nm & 57 & 47 & 12 & 12 & $12 \le 15$: keep \\
2470\,nm & 75 & 65 & 30 & 30 & $30 > 15$: drop \\
\bottomrule
\end{tabular}
\end{center}
(Script check\tagC.) A dropped band would meet wavelength networks that never saw anything there; dropping it is safer
than letting them extrapolate.
\end{worked}

The docstring states the trap in its second half: it ``also sets the wavelengths for the next predict\_full\_scene /
segment\_full\_scene''. Skip it and the model keeps the \emph{previous} scene's wavelengths - a mismatch that raises no
error, the same family of silent bug as \cref{ch:sortbug}. In the worker, the branch is taken only for
\code{band_grid == "native"}, and everything after it - reconstruction, scoring, segmentation, the reconstruction
preview - uses the kept cube and wavelengths. The stats record how many bands the model read (\texttt{bands\_read}).
On the PRISMA bench V2 reads 185 bands of set1 and 186 of set4.

\section{Score first, then zero the cloud}\label{sec:shipping:order}

\subsection{How the head turns a scene into classes}

\index{segmentation!tiled}
V2 has a segmentation head (\cref{ch:seghead}). \code{segment_full_scene} runs it over the scene in $128 \times 128$
windows at stride 64, averages the sigmoid probabilities where windows overlap, and assigns each valid pixel one class.
Three pieces of arithmetic hide in that sentence.

\paragraph{Padding so the windows fit.} With window $T = 128$ and stride $S = 64$, a side of length $n > T$ is padded
by $\text{pad} = (-(n - T)) \bmod S$, the smallest amount that makes $n + \text{pad} - T$ a multiple of $S$. For $n =
300$: $n - T = 172$; the next multiple of 64 at or above 172 is 192; $\text{pad} = 192 - 172 = 20$; padded side 320;
window starts $0, 64, 128, 192$ (since $320 - 128 = 192$). For $n = 256$: $256 - 128 = 128 = 2 \times 64$, pad 0,
starts $0, 64, 128$\tagC.

\paragraph{Coverage and averaging.} A row $y$ is covered by every window whose start $s$ satisfies $s \le y < s + 128$.
With starts $0, 64, 128, 192$, row 100 is covered by the windows starting at 0 and 64 - two of them; row 10 only by the
window at 0. In two dimensions an interior pixel is covered by $2 \times 2 = 4$ windows, and its probability for each
class is the mean of four predictions. If the four windows say cloud 0.62, 0.48, 0.55 and 0.39, the pixel's cloud
probability is $(0.62 + 0.48 + 0.55 + 0.39) / 4 = 2.04 / 4 = 0.51$.

\paragraph{Threshold and priority.} \index{class priority}
A class is ``on'' where its averaged probability exceeds 0.5. The code then assigns classes in a fixed order, each
assignment overwriting the last:

\begin{lstlisting}
cls = np.zeros((H, W), np.int64)                        # 0 land
cls[prob[0] > 0.5] = 1                                  # water
cls[prob[3] > 0.5] = 3                                  # shadow
cls[prob[4] > 0.5] = 4                                  # snow
cls[(np.maximum(prob[1], prob[2]) if thin_cloud else prob[1]) > 0.5] = 2   # cloud (or thin)
return cls * valid.astype(bool), prob
\end{lstlisting}

Because the last write wins (the same rule as \cref{sec:shipping:codename}, this time on an array), the effective
priority is cloud $>$ snow $>$ shadow $>$ water $>$ land. \code{CLASS_NAMES = ("land", "water", "cloud", "shadow",
"snow")} gives the class-map ids 0 to 4.

\begin{worked}[label={we:shipping:priority}]{One pixel, two cloud rules}
A pixel's averaged probabilities (water, cloud, thin cloud, shadow, snow) are $(0.70, 0.30, 0.60, 0.55, 0.10)$.
\begin{enumerate}
\item Water $0.70 > 0.5$: class 1.
\item Shadow $0.55 > 0.5$: overwrite, class 3.
\item Snow $0.10$: no change.
\item Cloud, \texttt{thin\_cloud=True}: $\max(0.30, 0.60) = 0.60 > 0.5$: overwrite, class 2 (cloud). Excluded.
\item Cloud, \texttt{thin\_cloud=False}: $0.30$, not above 0.5: stays class 3 (shadow). Not excluded.
\end{enumerate}
\end{worked}

\begin{sourcenote}[: thick cloud only, since 2026-10-29]
At the FT04 ship, ``cloud'' meant cloud or thin cloud. Since 2026-10-29 the resolver sets \code{thin_cloud_is_cloud=False}
for \texttt{spectral\_stem\_finetune} (\code{backend/allotrope/foundation_models/resolver.py}, the capabilities entry), and
the worker passes it to \code{segment_full_scene(..., thin_cloud=m.thin_cloud_is_cloud)}: V2 excludes \emph{thick} cloud
only. The reason is Indradhanu Heavy's thin-cloud output, which fires on haze (75.4\% of PRISMA set4's valid pixels and
209 of its 230 targets; \file{docs/lld/spectral-stem.md}). \texttt{stats.json} records the rule as \texttt{cloud\_rule}
(\texttt{"thick"} or \texttt{"cloud\_or\_thin"}). \Cref{ch:heavy} tells that story.
\end{sourcenote}

\subsection{Why zero after scoring}

The worker then zeros the model's \code{exclude_classes} - for V2, cloud:

\begin{lstlisting}
# _anomaly_scoring_run.py, after compute_score(...)
class_map, _ = inferencer.segment_full_scene(model_cube, spatial_valid.astype(bool),
                                             thin_cloud=m.thin_cloud_is_cloud)
excluded = np.isin(class_map, [inferencer.CLASS_NAMES.index(c) for c in m.exclude_classes])
score = np.where(excluded, 0.0, score).astype(np.float32)          # (H, W)
\end{lstlisting}

Why after, and not ``remove the cloud from the mask, then score''? Residual Mahalanobis (\cref{ch:residualmaha}) first
fits its whitening - the spread of the residual - over every pixel in the mask, and measures each pixel in units of that
spread. Which pixels are in the fit changes every score.

\begin{worked}[label={we:shipping:fitset}]{The same pixel, twice the score}
An illustration in one band, not a measurement\tagA: four land residuals $0.02, -0.02, 0.02, -0.02$, two cloud residuals
$0.06, -0.06$, and a target residual $0.08$. All have mean 0, so the variance is the mean of the squares.

\emph{Fit on all six.}
\begin{align*}
\text{squares:}\quad & 4 \times 0.0004 + 2 \times 0.0036 = 0.0016 + 0.0072 = 0.0088\\
\text{variance:}\quad & 0.0088 / 6 = 0.001467\\
\text{spread:}\quad & \sqrt{0.001467} = 0.0383\\
\text{target:}\quad & 0.08 / 0.0383 = 2.09 \text{ spreads}
\end{align*}
\emph{Fit on land only.}
\begin{align*}
\text{variance:}\quad & 0.0016 / 4 = 0.0004\\
\text{spread:}\quad & \sqrt{0.0004} = 0.0200\\
\text{target:}\quad & 0.08 / 0.0200 = 4.00 \text{ spreads}
\end{align*}
(Script: 2.089 and 4.000\tagC.)
\end{worked}

Same pixel, nearly twice the score. Neither order is wrong in principle, but \textbf{the benchmark measured ``fit on
everything valid, then zero the cloud''}, and the product must run what was measured. The evaluator's other order
(\code{fit_on_kept: true} in \code{app/evaluation/scorers.py}) was tried on FT03 and was no better: 15.7\% against 16.5\%
caught at the top 0.1\%\tagM{docs/lld/evaluation.md}.

A note on the keep mask: if the analyst also attaches an upstream \code{scene_segmentation} keep mask, that mask does feed
the fit, as it always has. The parity run below used no upstream keep mask - only the valid pixels - which is the
configuration the numbers describe.

\subsection{What is written}

Three things are written per model, under \file{models/<codename>/}:
\begin{center}
\small
\begin{tabular}{ll}
\toprule
output & contents \\
\midrule
\texttt{mask\_cloud.tif} & uint8, 1 where the score was zeroed \\
\texttt{class\_map.tif} & uint8, the head's classes, 0 land to 4 snow \\
\texttt{stats.json} keys & \texttt{bands\_read}, \texttt{excluded\_classes}, \texttt{cloud\_rule}, \texttt{excluded\_percent} \\
\bottomrule
\end{tabular}
\end{center}
\texttt{excluded\_percent} is \code{round(100 * excluded / valid, 3)}.

\begin{worked}[label={we:shipping:percent}]{excluded\_percent on a small tile}
A $4 \times 4$ tile has 14 valid pixels, of which 3 are cloud:
\begin{align*}
100 \times 3 / 14 &= 300 / 14 = 21.42857\ldots\\
\text{rounded to 3 places} &= 21.429
\end{align*}
\end{worked}

The keys are additive: the frontend reads stats by key, so new keys break nothing - and nothing in the viewer shows them
yet (\file{docs/09-known-issues.md}).

\begin{figure}[htbp]
\centering
\begin{tikzpicture}[node distance=5mm and 6mm, font=\small,
  box/.style={draw, rounded corners=2pt, align=center, minimum height=8mm, text width=3.4cm, fill=white, font=\footnotesize},
  arr/.style={-{Stealth[length=2mm]}, thick}]
\node[box, fill=spat!10] (pick) {analyst picks\\ ``Indradhanu V2''};
\node[box, right=of pick] (cat) {\texttt{list\_catalog}: manifest\\ with embedded config};
\node[box, right=of cat] (grid) {\texttt{\_check\_band\_grid}:\\ native, not 165};
\node[box, fill=spec!10, below=of grid] (read) {\texttt{read\_bands}: keep\\ bands within 15\,nm};
\node[box, left=of read] (pred) {\texttt{predict\_full\_scene}\\ (two checkerboards)};
\node[box, fill=loss!15, left=of pred] (score) {residual Mahalanobis\\ fit on valid pixels};
\node[box, fill=attn!10, below=of score] (seg) {\texttt{segment\_full\_scene}\\ 128 px, stride 64};
\node[box, fill=loss!15, right=of seg] (zero) {zero cloud\\ \emph{after} scoring};
\node[box, right=of zero] (out) {\texttt{anomaly\_score.tif},\\ \texttt{mask\_cloud.tif}, stats};
\draw[arr] (pick) -- (cat); \draw[arr] (cat) -- (grid); \draw[arr] (grid) -- (read);
\draw[arr] (read) -- (pred); \draw[arr] (pred) -- (score); \draw[arr] (score) -- (seg);
\draw[arr] (seg) -- (zero); \draw[arr] (zero) -- (out);
\end{tikzpicture}
\caption{The worker's path for Indradhanu V2, in order. Every box is one of the decisions of this chapter; the parity run
(\cref{sec:shipping:parity}) runs the same boxes outside the worker and compares with the evaluator.}
\label{fig:shipping:path}
\end{figure}

\section{The parity run: the guard}\label{sec:shipping:parity}

\index{parity run}
There are no behavioural backend tests in this repo. \code{tests/test_backend/test_native_band_models.py} checks the
plumbing: the resolver reads an embedded config, refuses a missing or mismatched one, applies the thick-cloud rule, and the
grid check accepts each grid for its own kind of model and refuses the other. It does not check numbers. The guard for the
numbers is \file{research/experiments/indradhanu-finetune/FT04-final-joint/product_parity.py}: it builds the inferencer
\emph{from the shipped manifest} through the resolver's \code{_from_manifest} and \code{get_inferencer}, then runs
\code{read_bands}, \code{predict_full_scene}, \code{compute_score(method="residual_mahalanobis")},
\code{segment_full_scene} and the zeroing, exactly as the worker does, and compares against the evaluator's FT04 numbers on
the PRISMA bench.

\begin{table}[htbp]
\centering
\caption{The parity run, 2026-10-19: product path against the evaluator, PRISMA bench.}
\label{tab:shipping:parity}
\small
\begin{tabular}{lcccc}
\toprule
scene & bands read & pAUC plain & pAUC cloud-excluded & excluded \\
\midrule
set1 & 185 & 0.2460 & -- & -- \\
set4 & 186 & 0.2106 & 0.2199 & 11.83\% of valid pixels \\
\bottomrule
\end{tabular}

\smallskip
{\footnotesize Source: \tagM{docs/lld/spectral-stem.md} (Serving), produced by
\file{research/experiments/indradhanu-finetune/FT04-final-joint/product_parity.py}.}
\end{table}

Identical to the evaluator's run; ROC-AUC within 0.0005. Two things made ``identical'' possible. The evaluator
(\code{app/evaluation/models.py}) now calls the same \code{read_bands} and \code{segment_full_scene}, so there is one
implementation of band reading and segmentation, not two that could drift. And \cref{ch:sortbug}'s lesson - a
re-implemented data path that quietly disagreed with the product's - is the reason a parity run exists at all.

These per-scene pAUCs are not comparable with the four-scene 0.170: they are single scenes. What matters is only that
product and evaluator agree. If set4 has 1,000,000 valid pixels, 11.83\% of them is $0.1183 \times 1{,}000{,}000 =
118{,}300$ pixels scored 0 by the exclusion.

\begin{drift}
A second parity run was made when Heavy took over the codename (2026-10-29, thick-cloud rule): pAUC set1 0.2617, set2
0.1874, set3 0.0861, set4 0.2217 (mean 0.1892), matching the evaluator to four decimals (\cref{we:heavy:parity}). CPU
cost inside the worker: on the image's 2 torch threads one PRISMA scene took 1,819\,s with Heavy and 1,058\,s with
FT04, mostly \code{predict_full_scene} at stride 32; since 2026-10-29 the compose file raises the worker to 16 threads
(\code{WORKER_THREADS}), and the same scene takes 454\,s and 274\,s, about four times faster
(\file{docs/09-known-issues.md}; \cref{we:heavy:threads}).
\end{drift}

\begin{keypoint}
Unit tests prove the plumbing; only a run of the real checkpoint on the real scenes, through the product's own path,
proves the numbers.
\end{keypoint}

\section{What did not change}\label{sec:shipping:unchanged}

\begin{itemize}
\item \textbf{Defaults.} Every hyperspectral default still picks ``Indradhanu'' (v0.2.0); \code{band_filter_apply} still
  defaults to \code{use_common_wavelength_grid: true}. Phase 2 would change both.
\item \textbf{EnMAP's band filter} still removes the provider's cloud, cloud-shadow and haze pixels
  (\code{quality_masks_to_apply}) before scoring, unlike the PRISMA evaluation behind V2's numbers. Kept as the product's
  convention (decision, 2026-10-19).
\item \textbf{The detection case.} Like for like, both with cloud exclusion on FT04's masks, the real bench reads v0.2.0
  0.165 against FT04 0.170 - a tie (\cref{ch:ftruns}). What V2 adds is the masks, native bands and better
  reconstruction, in one model.
\item \textbf{One product fix rode along.} \code{scene_segmentation}'s brightness took the mean of ``band 0 up to the VNIR
  edge''. On native PRISMA, whose bands run downwards (\cref{ch:sortbug}), band 0 is the far SWIR.
  \code{app/utils/spectral_indices/scene_indices.py} now picks bands by wavelength (\code{find_band}: the index of the
  nearest wavelength, \code{np.argmin(np.abs(wavelengths - target_nm))}), which works in any order.
\end{itemize}

Still open (\file{docs/09-known-issues.md}): the viewer has no panel for the masks, and the CPU cost on the stack, measured
since (drift box above), is several times v0.2.0's. Anything that runs the worker image outside compose is back on 2
threads, four times slower.

\begin{confusions}
\confusion{Indradhanu V2 replaced Indradhanu.}{No. It is opt-in. ``Indradhanu'' is still v0.2.0 and still the default.}
\confusion{The codename collision would only matter if someone picked the name.}{It matters for everyone: the defaults pick
``Indradhanu'' for you. A collision would have re-routed every default action to FT04 without a single config change.}
\confusion{The grid check counts bands.}{It compares wavelengths with the grid using \texttt{np.allclose}. 165 native bands
are still native.}
\confusion{Zeroing cloud before scoring is the same thing, earlier.}{No. It changes which pixels the whitening is fitted on,
so it changes every score, not only the cloud's.}
\confusion{mask\_cloud.tif and the upstream keep\_mask.tif are the same kind of mask.}{\texttt{keep\_mask.tif} comes from
\texttt{scene\_segmentation} (thresholds) and decides the fit. \texttt{mask\_cloud.tif} comes from V2's own head and is
applied after scoring. Both can be present.}
\confusion{The manifest's params is the number of trainable weights.}{It is every state-dict entry except the normalisation
and wavelength buffers, so it includes the BatchNorm's 65 running-statistics entries: 5,886,451 against 5,886,386 trainable.}
\confusion{The unit tests prove the product gives the benchmark's numbers.}{They prove the plumbing (config, grid, cloud
rule). The numbers are guarded by the parity run, which needs the checkpoint and the bench scenes and is not part of
pytest.}
\end{confusions}

\begin{checkyourself}
\item\label{ex:shipping:1} Suppose FT04's directory had been named \texttt{a\_spectral\_stem} and both manifests said
  ``Indradhanu''. Which model would the defaults route to, and why? Show the character codes.
\item\label{ex:shipping:2} A trained record contains 1385 and 1460\,nm and nothing between. Does \texttt{read\_bands} keep a
  band at 1400\,nm? At 1420\,nm? At 1445\,nm? Show the distances.
\item\label{ex:shipping:3} A band filter output lists wavelengths equal to the 165-band grid plus 0.003\,nm at every band.
  Does \texttt{\_check\_band\_grid} treat it as on the grid? What about plus 0.03\,nm? Use \eqref{eq:shipping:allclose} at
  460 and at 2450\,nm.
\item\label{ex:shipping:4} A 6 $\times$ 5 tile has 26 valid pixels; V2 calls 4 of them cloud. Compute
  \texttt{excluded\_percent} by hand. Then a scene of side $n = 200$: what padding does \texttt{segment\_full\_scene} add, and
  where do its windows start?
\item\label{ex:shipping:5} In \cref{we:shipping:fitset}, add a third cloud residual of $0.06$ (seven pixels in all, mean no
  longer exactly 0 - use the mean of the squares anyway, as an approximation). Recompute the target's score with the fit on
  all pixels. Does adding cloud to the fit raise or lower the target's score?
\item\label{ex:shipping:6} An analyst picks both ``Indradhanu'' and ``Indradhanu V2'' in one \texttt{anomaly\_scoring} action
  on a default band filter. What happens, and what should they do instead?
\end{checkyourself}
\endgroup

\nextchapter{ch:scoringgains}
